\begin{figure}[ht]
\centering
\begin{tikzpicture}[>=Stealth, font=\small, node distance=2.1cm]
  \node[draw, rounded corners, fill=blue!8, minimum width=3.8cm, minimum height=1cm] (E)
    {$E/\F_p$};
  \node[draw, rounded corners, fill=green!8, right=4.2cm of E, minimum width=3.8cm, minimum height=1cm] (Ep)
    {$E^{(p)}/\F_p$};
  \node[draw, rounded corners, fill=orange!10, below=1.9cm of $(E)!0.5!(Ep)$, minimum width=5.6cm, minimum height=1cm] (corr)
    {Hecke 对应 $T_p$: 枚举阶 $p$ 子群 $D\subset E[p]$};
  \node[draw, rounded corners, fill=purple!10, below left=of corr, minimum width=3.9cm, minimum height=1cm] (kerf)
    {$D=\ker(\phi_p)$};
  \node[draw, rounded corners, fill=red!8, below right=of corr, minimum width=3.9cm, minimum height=1cm] (kerv)
    {$D=\ker(V_p)$};

  \draw[->, very thick] (E) -- node[above] {$\phi_p:(x,y)\mapsto (x^p,y^p)$} (Ep);
  \draw[->, very thick] (Ep) -- node[below] {$V_p$} (E);
  \draw[->, thick] (corr) -- (kerf);
  \draw[->, thick] (corr) -- (kerv);
  \draw[dashed, ->, thick] (kerf) -- node[left] {Frobenius 方向} ($(E.south)!0.45!(Ep.south)$);
  \draw[dashed, ->, thick] (kerv) -- node[right] {Verschiebung 方向} ($(Ep.south)!0.45!(E.south)$);
\end{tikzpicture}
\caption{在特征 $p$ 中，Hecke 对应 $T_p$ 按阶 $p$ 子群分解，其中最自然的两块正是 Frobenius 核与 Verschiebung 核；这正导致 $T_p=\Frob_p+V_p$。}
\label{fig:ch08-frobenius-action}
\end{figure}
